\section{Flusso di Continuous Integration}
Il Listato \ref{lst:ci-workflow} descrive il flusso di Continuous Integration codificato tramite sintassi YAML per la piattaforma GitHub Actions. Il processo automatizza il checkout del codice, la risoluzione dei submodules privati e demanda a un'azione esterna la compilazione dell'immagine container, condizionando l'esecuzione alla presenza di un versionamento semantico formale del codice sorgente.

\begin{lstlisting}[label={lst:ci-workflow}]
name: Dev CI, docker image build and push

on:
  push:
    branches:
      - main
      - master

permissions:
  id-token: write
  contents: write
  pull-requests: write
  issues: write

env:
  IMAGE_CACHE_DIR: /tmp/cache/docker-image
  IMAGE_CACHE_KEY: cache-image
  IMAGE_KEY_VERSION: 1
  SERVICE_NAME: default
  REPO_NAME: ${GITHUB_REPOSITORY#$GITHUB_REPOSITORY_OWNER/}

jobs:
  build-and-push:
    name: Setup, Build and Publish
    runs-on: ubuntu-latest

    steps:
      - name: Checkout code
        uses: actions/checkout@v4
        with:
          submodules: false
          
      - uses: actions/create-github-app-token@v1
        id: app-token
        with:
          app-id: ${{ vars.COMMON_SUBLIBS_APP_ID }}
          private-key: ${{ secrets.COMMON_SUBLIBS_APP_PRIVATE_KEY }}
          owner: ${{ github.repository_owner }}
          
      - name: Checkout Submodule con Token Personalizzato
        run: |
          git config --global url."https://git:${{ steps.app-token.outputs.token }}@github.com/".insteadOf "https://github.com/"
          git config --global url."https://git:${{ steps.app-token.outputs.token }}@github.com/".insteadOf "git@github.com:"
          git submodule update --init --recursive
          
      - uses: devredops/cicd-release-action@main
        id: release
        
      - name: Set ENV variables
        run: |
          echo "REPO_NAME=${GITHUB_REPOSITORY#$GITHUB_REPOSITORY_OWNER/}" >> $GITHUB_ENV
          
      - name: Build and Push Docker Image to ECR
        if: ${{ steps.release.outputs.release-created }}
        uses: Company/gha-build-and-push-on-ecr@main
        with:
          aws-role-to-assume: ${{ vars.AWS_ROLE_TO_ASSUME_OLD }}
          ecr-repository: ${{ env.REPO_NAME }}
          release-version: ${{ steps.release.outputs.release-version }}
          
      - name: Build and Push Docker Image to ECR
        if: ${{ steps.release.outputs.release-created }}
        uses: Company/gha-build-and-push-on-ecr@main
        with:
          aws-role-to-assume: ${{ vars.AWS_ROLE_TO_ASSUME }}
          ecr-repository: ${{ env.REPO_NAME }}
          release-version: ${{ steps.release.outputs.release-version }}
\end{lstlisting}


\section{Azione Composita per la Pubblicazione su Amazon ECR}
Il Listato \ref{lst:ecr-composite-action} espone l'azione GitHub composita richiamata dalla pipeline principale. Il codice astrae la logica di autenticazione al provider cloud tramite ruoli IAM, gestisce l'istanziazione condizionale del registro contenitori qualora non sia presente, e orchestra l'effettiva compilazione e pubblicazione dell'immagine, garantendo la standardizzazione del processo di build aziendale.

\begin{lstlisting}[label={lst:ecr-composite-action}]
name: 'gha-build-and-push-on-ecr'
description: 'Automates the process of building, tagging, and pushing a Docker image to Amazon ECR.'

inputs:
  aws-role-to-assume:
    description: "AWS Role to assume"
    required: true
  aws-region:
    description: "AWS Region"
    default: "eu-west-1"
    required: true
  ecr-repository:
    description: 'Registry Domain'
    required: true
  build-platform:
    description: "Build Platform"
    default: "linux/arm64"
    required: true
  build-folder:
    default: "."
    description: "Folder where Dockerfile is located"
    required: true
  release-version:
    description: "Release version es: 1.0.0"
    required: true

outputs:
  image-ecr-url:
    description: ""
    value: "https://${{ inputs.ecr-repository }}:${{ inputs.release-version }}"

runs:
  using: "composite"
  steps:
    - name: Checkout
      uses: actions/checkout@v4
      
    - name: Configure AWS Credentials
      uses: aws-actions/configure-aws-credentials@v4
      with:
        role-to-assume: ${{ inputs.aws-role-to-assume }}
        aws-region: ${{ inputs.aws-region }}
        
    - name: Login to Amazon ECR
      id: login-ecr
      uses: aws-actions/amazon-ecr-login@v2
      
    - name: Create Amazon ECR, if not exist
      shell: bash
      run: |
        aws ecr describe-repositories --repository-names ${{ inputs.ecr-repository }} || aws ecr create-repository --image-tag-mutability IMMUTABLE --image-scanning-configuration scanOnPush=true --repository-name ${{ inputs.ecr-repository }}
        
    - name: Apply ECR Repository Policy
      shell: bash
      run: |
        CURRENT_POLICY=$(aws ecr get-repository-policy --repository-name ${{ inputs.ecr-repository }} --query 'policyText' --output text 2>/dev/null || echo "")
        LOCAL_POLICY=$(cat ${{ github.action_path }}/ecr-policy.json)

        if [ "$(echo "$CURRENT_POLICY" | jq -S .)" == "$(echo "$LOCAL_POLICY" | jq -S .)" ]; then
          echo "ECR repository policy is already up-to-date. Skipping."
        else
          echo "Applying ECR repository policy."
          aws ecr set-repository-policy --repository-name ${{ inputs.ecr-repository }} --policy-text file://${{ github.action_path }}/ecr-policy.json
        fi
        
    - name: Set up QEMU
      uses: docker/setup-qemu-action@v3
      
    - name: Build, tag, and push image to Amazon ECR
      id: build-publish
      shell: bash
      env:
        ECR_REPOSITORY: ${{ steps.login-ecr.outputs.registry }}/${{ inputs.ecr-repository }}
        IMAGE_TAG: ${{ inputs.release-version }}
      run: |
        docker build --platform ${{ inputs.build-platform }} --provenance=false -t "$ECR_REPOSITORY:$IMAGE_TAG" ${{ inputs.build-folder }}
        docker push "$ECR_REPOSITORY:$IMAGE_TAG"
        echo "IMAGE $IMAGE_TAG is pushed to $ECR_REPOSITORY"
        echo "image_tag=$IMAGE_TAG" 
        echo "full_image=$ECR_REPOSITORY:$IMAGE_TAG"
\end{lstlisting}


\section{Flusso di Continuous Deployment}
Il Listato \ref{lst:cd-workflow} dettaglia la pipeline di Continuous Deployment per la propagazione degli artefatti. Questa action richiede, a differenza delle precedenti, un innesco esplicito (evento \texttt{workflow\_dispatch}) da parte dell'operatore per selezionare l'ambiente di destinazione e distribuire le immagini validate verso gli account segregati di dev o prod.

\begin{lstlisting}[label={lst:cd-workflow}]
name: Deploy to Environment

on:
  workflow_dispatch:
    inputs:
      environment:
        description: 'The environment to deploy to'
        required: true
        type: choice
        options:
        - dev
        - prod
        default: 'dev'

permissions:
  id-token: write
  contents: read

env:
  IMAGE_CACHE_DIR: /tmp/cache/docker-image
  IMAGE_CACHE_KEY: cache-image
  IMAGE_KEY_VERSION: 1
  SERVICE_NAME: default
  REPO_NAME: ${GITHUB_REPOSITORY#$GITHUB_REPOSITORY_OWNER/}

jobs:
  deploy:
    name: Deploy to ${{ github.event.inputs.environment }} Environment
    runs-on: ubuntu-latest

    steps:
      - name: Set release version tag in env
        run: echo "RELEASE_VERSION=${GITHUB_REF#refs/*/v}" >> $GITHUB_ENV

      - name: Echo
        run: |
          echo "Tag founded is: ${{ env.RELEASE_VERSION }}"
          echo "repo name is: ${{ env.REPO_NAME }}"

      - name: Get Release Tag
        run: echo "RELEASE_TAG=${{ github.ref_name }}" >> $GITHUB_ENV

      - name: Echo Release Tag
        run: "echo Deploying release tag: ${{ env.RELEASE_TAG }} to ${{ github.event.inputs.environment }}"

      - name: Set Lambda Function Name
        run: |
            REPO_NAME=${{ env.REPO_NAME }}
            LAMBDA_NAME="${{ env.REPO_NAME }}-${{ github.event.inputs.environment }}"
            echo "LAMBDA_FUNCTION_NAME=$LAMBDA_NAME" >> $GITHUB_ENV

      - name: Update Lambda Image (Dev)
        if: github.event.inputs.environment == 'dev'
        uses: Company/gha-lambda-image-upload@main
        with:
          lambda-function-name: ${{ env.LAMBDA_FUNCTION_NAME }}
          ecr-repository-uri: ${{ vars.ECR_REPOSITORY_URI }}${{ env.REPO_NAME }}
          image-tag: ${{ env.RELEASE_VERSION }}
          aws-role-to-assume: ${{ vars.AWS_ROLE_DEPLOY_DEV }}
          aws-region: ${{ vars.AWS_REGION }}

      - name: Update Lambda Image (Prod)
        if: github.event.inputs.environment == 'prod'
        uses: Company/gha-lambda-image-upload@main
        with:
          lambda-function-name: ${{ env.LAMBDA_FUNCTION_NAME }}
          ecr-repository-uri: ${{ vars.ECR_REPOSITORY_URI }}${{ env.REPO_NAME }}
          image-tag: ${{ env.RELEASE_VERSION }}
          aws-role-to-assume: ${{ vars.AWS_ROLE_DEPLOY_PROD }}
          aws-region: ${{ vars.AWS_REGION }}
\end{lstlisting}

\section{Azione Composita per l'aggiornamento dell'immagine ECR}
Il Listato \ref{lst:lambda-update-action} espone l'azione GitHub composita responsabile dell'aggiornamento del codice applicativo. Il modulo automatizza l'invocazione delle interfacce a riga di comando (CLI) del provider cloud, istruendo l'infrastruttura ad agganciare il nuovo identificativo di versione dell'immagine container e implementando un meccanismo di attesa sincrona per validare l'effettivo completamento del rilascio.

\begin{lstlisting}[label={lst:lambda-update-action}]
name: 'Update Lambda ECR Image'
description: 'Updates an AWS Lambda function to use a new ECR image tag.'

inputs:
  lambda-function-name:
    description: "The name of the Lambda function to update."
    required: true
  ecr-repository-uri:
    description: "The full URI of the ECR repository."
    required: true
  image-tag:
    description: "The new image tag to deploy (e.g., the release version)."
    required: true
  aws-role-to-assume:
    description: "The ARN of the AWS IAM role to assume for deployment."
    required: true
  aws-region:
    description: "The AWS region where the resources are located."
    required: true
    default: "eu-west-1"

runs:
  using: "composite"
  steps:
    - name: Configure AWS Credentials
      uses: aws-actions/configure-aws-credentials@v4
      with:
        role-to-assume: ${{ inputs.aws-role-to-assume }}
        aws-region: ${{ inputs.aws-region }}

    - name: Log Update Details
      shell: bash
      run: |
        echo "Updating Lambda function: ${{ inputs.lambda-function-name }}"
        echo "Using image URI: ${{ inputs.ecr-repository-uri }}:${{ inputs.image-tag }}"

    - name: Update Lambda Function Code
      shell: bash
      run: |
        aws lambda update-function-code \
          --function-name ${{ inputs.lambda-function-name }} \
          --image-uri ${{ inputs.ecr-repository-uri }}:${{ inputs.image-tag }}

    - name: Wait for Lambda Update to Complete
      shell: bash
      run: |
        echo "Waiting for the update to complete..."
        aws lambda wait function-updated-v2 --function-name ${{ inputs.lambda-function-name }}
        echo "Update for '${{ inputs.lambda-function-name }}' completed successfully."
\end{lstlisting}